%
% ring.tex -- ring
%
% (c) 2021 Prof Dr Andreas Müller, OST Ostschweizer Fachhochschule
%
\documentclass[tikz]{standalone}
\usepackage{amsmath}
\usepackage{times}
\usepackage{txfonts}
\usepackage{pgfplots}
\usepackage{csvsimple}
\usetikzlibrary{arrows,intersections,math,calc}
\definecolor{darkred}{rgb}{0.8,0,0}
\definecolor{darkgreen}{rgb}{0,0.6,0}
\begin{document}
\def\skala{1}
\def\rone{0.8}
\def\rtwo{2.5}
\pgfmathparse{0.5*(\rone+\rtwo)}
\xdef\rmiddle{\pgfmathresult}
\def\w{40}
\def\pfeilkopf#1{
	\draw[->,color=darkred,line width=1.2pt]
		($(C)+({#1+4}:\rtwo)+({#1+4-90}:0.1)$)
		--
		($(C)+({#1+4}:\rtwo)$);
}
\def\Pfeilkopf#1{
	\draw[->,color=darkred,line width=1.2pt]
		($(C)+({#1-6}:\rone)+({#1-6+90}:0.1)$)
		--
		($(C)+({#1-6}:\rone)$);
}
\def\Gammakopf#1{
	\draw[->,color=darkgreen,line width=1.2pt]
		({X(#1)},{Y(#1)})
		--
		({X(#1+0.1)},{Y(#1+0.1)});
}
\begin{tikzpicture}[>=latex,thick,scale=\skala,
declare function = {
	R(\t) = \rmiddle - 0.5*cos(\t-\w) + 0.2 * sin(5*\t);
	X(\t) = 5 + R(\t) * cos(\t);
	Y(\t) = 4 + R(\t) * sin(\t);
}]

\coordinate (P1) at (2,1);
\coordinate (P2) at (8,1);
\coordinate (P21) at (10,3);
\coordinate (P3) at (9,3);
\coordinate (P4) at (9,6);
\coordinate (P5) at (5,8);
\coordinate (P6) at (2,7);
\coordinate (P7) at (0.5,5);

\fill[color=gray!20]
	(P1) to[out=-10,in=-135]
	(P2) to[out=45,in=-90]
	(P21) to[out=90,in=45]
	(P3) to[out=-135,in=180]
	(P4) to[out=0,in=0]
	(P5) to[out=180,in=100]
	(P6) to[out=-80,in=140]
	(P7) to[out=-40,in=170]
	cycle;

%\fill (P1) circle[radius=0.05];
%\fill (P2) circle[radius=0.05];
%\fill (P21) circle[radius=0.05];
%\fill (P3) circle[radius=0.05];
%\fill (P4) circle[radius=0.05];
%\fill (P5) circle[radius=0.05];
%\fill (P6) circle[radius=0.05];
%\fill (P7) circle[radius=0.05];

\coordinate (C) at (5,4);
\coordinate (D) at ($(C)+(\w:\rmiddle)$);

\fill[color=darkred!20] (C) circle[radius=\rtwo];

\begin{scope}
\clip (C) circle[radius=\rone];
\fill[color=gray!20]
	(P1) to[out=-10,in=-135]
	(P2) to[out=45,in=-90]
	(P21) to[out=90,in=45]
	(P3) to[out=-135,in=180]
	(P4) to[out=0,in=0]
	(P5) to[out=180,in=100]
	(P6) to[out=-80,in=140]
	(P7) to[out=-40,in=170]
	cycle;
\end{scope}

\begin{scope}[xshift=5cm,yshift=4cm]
	\fill[color=white] (0,0.4) [rotate=-30] ellipse(0.5cm and 0.3cm);
\end{scope}

\draw[->,line width=0.5pt] (C) -- ++(-148:\rtwo);
\node at ($(C)+(-155:{0.65*\rtwo})$) {$r_2$};

\draw[->,line width=0.5pt] (C) -- ++(150:\rone);
\node at ($(C)+(165:{0.7*\rone})$) {$r_1$};


\draw[color=darkred,line width=1.2pt]
	(C) circle[radius=\rone];

\pfeilkopf{90}
\pfeilkopf{135}
\pfeilkopf{180}
\pfeilkopf{225}
\pfeilkopf{270}
\pfeilkopf{315}
\pfeilkopf{0}

\Pfeilkopf{135}
\Pfeilkopf{225}
\Pfeilkopf{315}

\draw[color=darkred,line width=1.2pt]
	(C) circle[radius=\rtwo];

\draw[color=blue,line width=1.8pt]
	($(C)+(\w:\rone)$) -- ($(C)+(\w:\rtwo)$);

\fill[color=gray!20]
	($(C)+(\w:\rmiddle)$) circle[radius=0.18];
\draw[color=blue,line width=0.6pt]
	($(C)+(\w:\rmiddle)$) circle[radius=0.18];
\draw[color=gray!20,line width=0.8pt]
	($(C)+(\w:{\rone-0.04})$) -- ($(C)+(\w:{\rtwo+0.04})$);

\draw[color=blue,line width=0.2pt] ($(D)+(135:0.18)$) -- (D);
\draw[->,color=blue,line width=0.2pt] ($(D)+(135:0.5)$) -- ($(D)+(135:0.18)$);
\node[color=blue] at ($(D)+(135:0.18)$) [above] {$\varepsilon$};

\fill[color=white] (D) circle[radius=0.05];
\draw[color=blue] (D) circle[radius=0.05];
\node[color=blue] at ($(D)+(0,0)$) [below right] {$a$};

\fill[color=white] (C) circle[radius=0.05];
\draw (C) circle[radius=0.05];
\node at (C) [below] {$z_0$};

\node[color=darkred] at ($(C)+(270:\rone)+(0,0.05)$) [below] {$\gamma_1^{-1}$};
\node[color=darkred] at ($(C)+(110:\rtwo)+(0,-0.1)$) [above left] {$\gamma_2$};

\draw[->] (-0.1,0) -- (10.3,0) coordinate[label={$\operatorname{Re}$}];
\draw[->] (0,-0.1) -- (0,8.3) coordinate[label={right:$\operatorname{Im}$}];
\node at (10,8) {$\mathbb{C}$};

\node at (2,2) {$U$};

\node at ($(C)+(\w:\rtwo)$) [above right] {$A$};
\node at ($(C)+(\w:\rone)$) [below left] {$B$};

\draw[color=darkgreen,line width=1.2pt]
	plot[domain=0:360,samples=360] ({X(\x)},{Y(\x)});

\foreach \x in {5,41,...,360}{
	\Gammakopf{\x}
}
\node[color=darkgreen] at ({X(125)},{Y(125)}) [above] {$\gamma$};

\end{tikzpicture}
\end{document}

